%!TEX root = ../template.tex
%%%%%%%%%%%%%%%%%%%%%%%%%%%%%%%%%%%%%%%%%%%%%%%%%%%%%%%%%%%%%%%%%%%%
%% abstrac-en.tex
%% NOVA thesis document file
%%
%% Abstract in English
%%%%%%%%%%%%%%%%%%%%%%%%%%%%%%%%%%%%%%%%%%%%%%%%%%%%%%%%%%%%%%%%%%%%

\typeout{NT FILE abstrac-en.tex}%

Modern online services serve hundreds of millions of users from data centers
spread across the globe, and must reconcile three conflicting requirements:
low latency, high availability, and data consistency. As the number of data
centers grows, replicating every object everywhere becomes prohibitive in
storage, network and processing costs. \emph{Partial replication}, where each
data center replicates only a subset of the data, is a natural way to contain
these costs but introduces significant challenges: maintaining application
invariants under weak consistency, executing transactions whose objects are
not all replicated locally, and answering recurring analytical queries that
depend on data scattered across many data centers.

This thesis investigates how to design a geo-distributed database that
supports partial replication efficiently while still offering useful
correctness guarantees and a rich, expressive interface to applications. We
propose PotionDB, a geo-distributed database whose contributions rest on
three pillars: (i) a replication algorithm that combines partial,
non-uniform and geo-replication, propagating to each server only the
operations relevant to its locally replicated data; (ii) a transactional
algorithm providing PSI+cset consistency, with snapshot isolation within a
server and causal consistency across servers, supporting transactions whose
state may not be fully present at any single replica; and (iii) materialized
views maintained incrementally as CRDT updates, kept atomically consistent
with their underlying objects, and able to summarize data that is partitioned
across many partial replicas. PotionDB also includes a state-reconstruction
algorithm for old object versions based on logs of operations and effects,
reducing memory footprint compared to keeping full version histories.

Future work focuses on dynamic partitioning, per-transaction consistency
levels (including invariants on objects and views), an SQL-like front-end,
and a refined garbage-collection algorithm. The proposed system is
implemented as a working prototype and evaluated on Grid'5000 using
TPC-H, YCSB and microbenchmarks, showing that complex analytical queries
over partially replicated data can be answered with latency comparable to
basic object reads, while preserving the consistency between views and
their referred objects.

% Palavras-chave do resumo em Inglês
\begin{keywords}
Geo-distributed databases, Partial replication, Non-uniform replication,
Materialized views, CRDTs, Causal consistency, Snapshot isolation,
Transactions, Invariants
\end{keywords}
